\chapter{\MakeUppercase{Results and Discussion}}

\section{Model comparison}

I benchmarked each model against the 97-question evaluation set using the v2 knowledge base. Tables 7.1 and 7.2 report the quality metrics and latency measurements for each model evaluated on the 97-question test set.

\begin{table}[htbp]
\centering
\caption{Table 7.1: Model quality comparison (n=97, KB v2, 95\% bootstrap CIs in brackets)}
\label{tab:model_quality_comparison}
\renewcommand{\arraystretch}{1.3}
\resizebox{\textwidth}{!}{
\begin{tabular}{|l|l|l|l|}
\hline
\textbf{Model} & \textbf{KW Coverage} & \textbf{ROUGE-L} & \textbf{BERTScore F1} \\
\hline
TinyLlama-1.1B (base) & 0.3845 [0.32, 0.46] & 0.2076 [0.17, 0.25] & 0.1398 [0.08, 0.20] \\
TinyLlama + QLoRA v2  & 0.4634 [0.39, 0.53] & 0.2374 [0.20, 0.28] & 0.179 [0.12, 0.23]  \\
BioMistral-7B (\ac{gguf} Q4\_K\_M) & 0.3088 [0.23, 0.39] & 0.2988 [0.23, 0.37] & 0.2676 [0.19, 0.34] \\
\hline
\end{tabular}}
\end{table}

\begin{table}[htbp]
\centering
\caption{Model latency performance (warm cache, CPU-only environment)}
\label{tab:model_latency_comparison}
\renewcommand{\arraystretch}{1.3}
\begin{tabular}{|l|l|}
\hline
\textbf{Model} & \textbf{Latency (warm, ms)} \\
\hline
TinyLlama-1.1B (base) & 115,451 \\
TinyLlama + QLoRA v2  & 35,282  \\
BioMistral-7B (\ac{gguf} Q4\_K\_M) & 143,149 \\
\hline
\end{tabular}
\end{table}

\begin{figure}[!ht]
    \centering
    \includegraphics[width=0.95\linewidth,height=0.85\textheight,keepaspectratio]{images/fig1_model_comparison.png}
    \caption{Metric comparison across generation backends. QLoRA fine-tuning raises keyword coverage the most; BioMistral leads on ROUGE-L and BERTScore F1.}
    \label{fig:model_comparison}
\end{figure}

QLoRA fine-tuning gave us the highest keyword coverage at 0.4634, an 8\% jump over the base model. Training on MedMCQA and PubMedQA dialogues helped the model pick up medical vocabulary. Interestingly, the BioMistral model led on ROUGE-L and BERTScore but lagged on keyword coverage. This reflection of its style: where TinyLlama gives a punchy 2-3 sentence answer, BioMistral provides 4-6 sentences with more context. This extra prose dilutes the keyword density, lowering the coverage score even though the actual answer might be more helpful.

I found that the choice of model really depends on who is using it. A clinician might want QLoRA’s precision with terminology, but a patient would likely prefer the fuller, more discursive explanations that BioMistral produces. The latency for these models on a CPU-only setup is a real trade-off. BioMistral’s first response can take over three minutes (228 seconds) to load and run. This cost is worth it for quality, but it definitely pushes the system into the realm of ``asynchronous'' rather than ``real-time.''

\section{Ablation study}

Table 7.3 shows how removing each confidence signal affects mean confidence and keyword coverage across the 97 questions.

\begin{table}[htbp]
\centering
\caption{Ablation study on confidence signals (n=97)}
\label{tab:ablation}
\renewcommand{\arraystretch}{1.3}
\begin{tabular}{|l|l|l|}
\hline
\textbf{Configuration} & \textbf{Mean confidence} & \textbf{Keyword coverage} \\
\hline
All signals (full system) & 0.0161 & 0.3881 \\
Ablate retrieval signal & 0.0070 & 0.4254 \\
Ablate generation signal & 0.0036 & 0.4304 \\
Ablate consistency signal & 0.0284 & 0.3794 \\
Ablate source agreement & 0.1743 & 0.4493 \\
Ablate entity coverage & 0.0169 & 0.3811 \\
\hline
\end{tabular}
\end{table}

\begin{figure}[!ht]
    \centering
    \includegraphics[width=0.95\linewidth,height=0.85\textheight,keepaspectratio]{images/fig4_ablation_eval.png}
    \caption{Figure 7.2: Ablation impact on keyword coverage. Each bar shows the coverage change when the corresponding signal is removed. Source agreement has the largest effect.}
    \label{fig:ablation}
\end{figure}

Removing the source agreement signal produces the largest jump in both keyword coverage (+0.061) and mean confidence (+0.158). Source agreement is the signal that most actively penalises responses when retrieved sources conflict. Without it, the system does not hedge, so more specific answers pass through the confidence filter unmodified. Whether that is better depends on whether the source conflict was genuine disagreement or merely different angles on the same condition.

Removing the retrieval and generation signals each raise keyword coverage by roughly 0.037--0.042, which suggests both signals sometimes suppress responses that would have scored well on keywords. The entity coverage signal has the smallest ablation impact; answers that score well on the other signals tend to address the question's medical entities anyway, so this signal rarely changes the outcome.

Tables 7.4 and 7.5 compare the full system against three structural variants: dense-only retrieval (no BM25 or RRF), no XAI (confidence scoring removed), and no RAG (LLM generating purely from parametric memory).

\begin{table}[htbp]
\centering
\caption{System variant quality comparison (n=97, TinyLlama, 95\% bootstrap CIs in brackets)}
\label{tab:baseline_quality}
\renewcommand{\arraystretch}{1.3}
\resizebox{\textwidth}{!}{
\begin{tabular}{|l|l|l|l|}
\hline
\textbf{Variant} & \textbf{KW Coverage} & \textbf{ROUGE-L} & \textbf{BERTScore F1} \\
\hline
Full system (TinyLlama) & 0.3845 [0.32, 0.46] & 0.2076 [0.17, 0.25] & 0.1398 [0.08, 0.20] \\
Dense-only retrieval    & 0.4189 [0.35, 0.49] & 0.2335 [0.19, 0.28] & 0.1836 [0.12, 0.24] \\
No XAI                  & 0.4337 [0.37, 0.50] & 0.2263 [0.19, 0.27] & 0.1757 [0.12, 0.23] \\
No RAG (LLM only)       & 0.2364 [0.19, 0.29] & 0.1505 [0.13, 0.17] & 0.1142 [0.07, 0.16] \\
\hline
\end{tabular}}
\end{table}

\begin{table}[htbp]
\centering
\caption{A comparison of latency profiles between system variants utilizing TinyLlama in a strict CPU-bound environment}
\label{tab:baseline_latency}
\renewcommand{\arraystretch}{1.3}
\begin{tabular}{|l|l|l|}
\hline
\textbf{Variant} & \textbf{Latency (ms)} & \textbf{P90 (ms)} \\
\hline
Full system (TinyLlama) & 115,451 & 276,838 \\
Dense-only retrieval    & 35,605  & 56,631  \\
No XAI                  & 34,043  & 54,882  \\
No RAG (LLM only)       & 12,290  & 14,160  \\
\hline
\end{tabular}
\end{table}

When I removed retrieval entirely, keyword coverage dropped by nearly 15 points. This is solid proof that the knowledge base is doing the heavy lifting, not just taking up time. The dense-only variant actually outperformed the full system on all quality metrics, which caught me by surprise. The full system includes BM25, corrective RAG, and confidence scoring --- all of which add complexity and latency. Dense-only is faster and less filtered, so for raw scores on this specific test set, it wins. However, the value of the ``extra'' machinery isn't just in raw overlap; it's about reducing hallucinations and giving users an auditable trail through source attribution and confidence signals. Those are safety features that ROUGE-L can't measure.

\section{Knowledge base version comparison}

\begin{table}[htbp]
\centering
\caption{Comparing dense retrieval performance directly between knowledge base variants v1 and v2 on the 97-question evaluation subset}
\renewcommand{\arraystretch}{1.3}
\begin{tabular}{|l|l|}
\hline
\textbf{Metric} & \textbf{Value} \\
\hline
KB v1 keyword coverage & 0.3845 \\
KB v2 keyword coverage & 0.6765 \\
Delta ($\ac{delta}$) & +0.2920 \\
Decision threshold & 0.03 \\
Decision & Use KB v2 \\
KB v2 document count & 505,584 \\
Chunking method & RecursiveSentenceChunker \\
\hline
\end{tabular}
\end{table}

The 0.292 gain from switching to v2 is the largest single improvement in the entire project, larger than the QLoRA gain by a factor of nearly four. The gain comes from two places. First, the \texttt{RecursiveSentenceChunker} preserves sentence boundaries, so retrieved chunks contain complete medical statements rather than truncated sentences that BM25 would struggle to rank. Second, bulking up the document database (reaching 505,000 artifacts) gave the retriever a larger surface area to match specific edge-case symptoms.

\section{Calibration results}

\begin{table}[htbp]
\centering
\caption{Confidence calibration results (n=97)}
\renewcommand{\arraystretch}{1.3}
\begin{tabular}{|l|l|}
\hline
\textbf{Metric} & \textbf{Value} \\
\hline
ECE (raw confidence) & 0.1438 \\
Mean predicted confidence & 0.0245 \\
Mean empirical label & 0.433 \\
\hline
\end{tabular}
\end{table}

\begin{figure}[!ht]
    \centering
    \includegraphics[width=0.95\linewidth,height=0.85\textheight,keepaspectratio]{images/fig3_reliability_diagram.png}
    \caption{Reliability diagram for raw confidence scores. If the system was perfectly calibrated, its points would follow the diagonal line. Instead, the model consistently undervalues its own empirical accuracy.}
    \label{fig:reliability}
\end{figure}

When I calculated the ECE at 0.1438, I immediately saw that the raw scores fail completely to align with empirical accuracy. The wide gap between mean confidence (0.0245) and the average empirical label (0.433) highlights my decision to intentionally engineer an under-confident system. Because my source agreement signal deliberately tanks confidence whenever retrieved documents diverge, and since medical literature often safely explores conditions from different angles regardless, this conservative bias heavily compounds. I consider this extreme conservatism completely justified for medical informatics, though it enforces my requirement to apply Platt scaling to calibrate raw outputs into actionable probabilities.

After applying Platt scaling, a separate ECE figure is not reported because the Platt parameters were fitted on the same 97-question set used for evaluation. A proper hold-out calibration set would be needed for a reliable post-scaling ECE estimate.

\section{Latency analysis}

\begin{table}[htbp]
\centering
\caption{Per-stage latency breakdown (TinyLlama, CPU-only)}
\renewcommand{\arraystretch}{1.3}
\begin{tabular}{|l|l|}
\hline
\textbf{Stage} & \textbf{Mean latency (ms)} \\
\hline
Query enhancement & 0.02 \\
Retrieval (BM25 + dense + RRF) & 42,624 \\
Corrective RAG & 1,218 \\
Generation & 69,868 \\
Hallucination detection & 1,736 \\
Confidence scoring & 5 \\
Total (cold, p90) & 276,838 \\
Total (warm, mean) & 115,451 \\
\hline
\end{tabular}
\end{table}

\begin{figure}[!ht]
    \centering
    \includegraphics[width=0.95\linewidth,height=0.85\textheight,keepaspectratio]{images/fig5_latency.png}
    \caption{Per-stage latency breakdown. The 115-second warm path comes almost entirely from the retrieval (42.6s) and generation (69.9s) steps. The remaining filters and checks combine for under three seconds flat.}
    \label{fig:latency}
\end{figure}

When I tracked the latency profiles, I found that the retrieval block (42.6 seconds) and the language model generation block (69.9 seconds) completely devoured the 115-second warm-path duration. Every other component basically registers as a rounding error. My query enhancement scripts tick over in 0.02 milliseconds, and even the heavy DeBERTa hallucination checking pass only burns 1.7 seconds. The cold-path metric spikes worse at the 90th-percentile (hitting 277 seconds) simply because the very first query takes the hit of unpacking the massive BM25 index from disk into active memory.

We estimate that on modern GPU hardware, dense embedding speed would roughly halve the retrieval stage (pulling it toward 21.3 seconds) and TinyLlama generation would drop from 70 seconds to under 2 seconds. Given this, total inference latency remains explicitly lower-bounded by the ~21-second retrieval overhead unless further optimizations inside the RRF and BM25 index matching routines are performed.

\section{RAGAS evaluation}

RAGAS~\cite{es2023ragas} was used for supplementary evaluation using non-\ac{llm} metrics to avoid requiring an evaluator API key. On a 10-question sample:

\begin{table}[htbp]
\centering
\caption{RAGAS metrics (TinyLlama, non-LLM metrics only, two sample sizes)}
\renewcommand{\arraystretch}{1.3}
\begin{tabular}{|l|l|l|}
\hline
\textbf{Metric} & \textbf{n=8} & \textbf{n=10} \\
\hline
Context precision (non-LLM) & 0.125 & 0.100 \\
Context recall (non-LLM) & 0.143 & 0.100 \\
ROUGE score (F-measure) & 0.063 & 0.226 \\
Exact match & 0.000 & 0.000 \\
\hline
\end{tabular}
\end{table}

I noticed a significant shift in results between the $n=8$ and $n=10$ runs, which just shows how sensitive these small-sample metrics can be to the specific questions chosen. The low precision and recall scores from RAGAS seem to stem from a mismatch between the tool's non-\ac{llm} logic and my actual pipeline. Since the non-\ac{llm} version of context recall relies on exact token matching, it doesn't quite work for a system like mine where answers are paraphrased rather than copied verbatim from the sources.

\section{Discussion}

I found that the quality of the knowledge base matters far more than model tuning at this scale. Improving the chunking logic (v2) gave me a gain four times larger than I achieved through QLoRA fine-tuning. If this data shows one thing, it's that getting the processing right is more important than picking the perfect model.

Applying these explainability signals also forced a trade-off between coverage and honesty. While removing the source agreement signal made the coverage scores look better on paper, it meant the system stopped flagging cases where sources actually disagreed. In a medical context, I believe it's much safer for the system to admit when it's unsure rather than risk giving a confidently wrong or conflicting answer.

Finally, BioMistral is the standout for quality if you can handle the latency. It's about 35\% slower than TinyLlama on a \ac{cpu}, but the depth of its answers is noticeably better. I designed the architecture to allow swapping these backends with a simple parameter change, so users can prioritize quality or speed depending on their hardware. 

While the retrieval hit rate of 97.2\% is encouraging, the one miss I encountered (involving a rare surgical technique) shows that there are still gaps in specialized literature. The system handles clinical conversations and research abstracts well, but very niche surgical knowledge is still underrepresented.
